% Transcription par lecture visuelle directe, sans OCR.
% Statut : transcrit ; relecture indépendante à effectuer.
% Chapitre complet : p. 166 (fin)–186 (début).
% Source : PDF original, vue 083, côté gauche, page manuscrite 166, après le chapitre cinquième.
\chapter{Seconde suite de la préfecture de M. Chaulet}
\begin{center}
Passage du pape à Lyon et ouverture de Fourvières
\end{center}

\section*{Motifs du voyage du pape}
% Mention marginale de la page manuscrite 166.
La Congrégation suivoit modestement le sentier que sembloit lui avoir tracé la Providence ; elle désiroit vivement saisir toutes les occasions d'assurer à la religion quelque nouveau triomphe, c'étoit la pensée principale qui avoit présidé à sa fondation, elle remplissoit les cœurs de tous les congréganistes ; mais il ne leur étoit pas donné de pressentir d'avance quels moyens leur seroient offerts pour servir cette belle cause. Le souverain maître ne demandoit d'eux qu'en serviteurs fidèles, ils se maintinssent toujours dans la ferme disposition de ne rien négliger pour
% Source : PDF original, vue 083, côté droit, page manuscrite 167.
procurer sa gloire toutes les fois que l'occasion s'en présenteroit, mais il se réservoit de manifester successivement par les événements, les circonstances particulières où il daigneroit les employer.

C'est ainsi que nous avons déjà vu la bonté divine rendre la Congrégation un instrument de sa miséricorde envers un grand nombre de chrétiens, en se servant de son intervention pour les rassembler autour de la divine eucharistie, et les réunir dans la pensée de devenir de fervents adorateurs de N. S. J. C. dans le sacrement de son amour. Aujourd'hui une nouvelle mission lui est encore réservée ; quelque chétive, pauvre et misérable qu'elle soit encore, c'est elle qui va être chargée au moment de l'arrivée du souverain pontife dans notre patrie, de donner l'élan à la ville de Lyon, et par elle à toute la France, afin que le chef de l'Église soit reçu partout avec cette effusion de sentiments, qui en consolant les cœurs fidèles, ranimera la foi des foibles. Vous savez, Messieurs, à quelle occasion fut entrepris le voyage du souverain pontife. Napoléon venoit de se faire déclarer empereur ; dans l'espoir de consolider son usurpation, il appelle la religion à son secours ; il cherche à se grandir en la prenant pour son piédestal. Il faudra donc qu'elle vienne l'entourer avec toutes ses pompes. Le sacre dans la métropole de Rheims, ainsi qu'il avoit lieu pour nos rois, ne suffiroit pas à son ambition, ce sera des mains révérées du chef de l'Église qu'il exigera que l'onction sainte découle sur son front, comme autrefois sur celui de Charlemagne ; car comme cet empereur si célèbre dans nos fastes, il aspire à devenir le fondateur d'une nouvelle dynastie. En conséquence le pape reçoit en forme d'invitation, un ordre de se rendre à Paris pour cette cérémonie. Pie VII, pontife
% Source : PDF original, vue 084, côté gauche, page manuscrite 168.
rempli de vertus et d'abnégation de sa propre personne, est décidé d'avance à tous les sacrifices que sa conscience ne réprouvera pas ; il sait qu'il importe de ménager cet homme orgueilleux, qui d'un seul mot pourroit renouveller la persécution ; il se résigne donc et entreprend aux approches de l'hiver un voyage qui pouvoit lui paroître d'autant plus pénible qu'il étoit loin de s'attendre, ainsi qu'il l'a dit lui-même, de trouver encore tant de foi, dans notre malheureuse France.

Vous savez, Messieurs, que je prends la liberté d'entremêler à mon récit les réflexions que je pense qui pourront vous être utiles. Permettez-moi donc de vous faire remarquer ici cette main puissante de notre Dieu qui dirige toutes les volontés, et fait que les vaines délibérations que l'orgueil suggère à l'esprit de l'homme, servent à l'accomplissement des desseins de sa souveraine sagesse. Ainsi, on vit autrefois l'empereur Auguste ordonner le dénombrement de tous ceux qui étoient soumis à son empire ; sa vanité se complaisoit à contempler par là toute l'étendue de sa puissance ; mais le Très-Haut, en le laissant agir ainsi conduisoit la Ste Vierge à Bethléem où le Messie devoit naître, conformément aux prophéties. Aujourd'hui l'ambition orgueilleuse de celui qui s'est emparé du pouvoir est le motif apparent qui conduit en France le souverain pontife. Mais Dieu dans sa bonté, a d'autres desseins. Cette visite du vicaire de J. C. a pour but de répandre un baume salutaire sur les plaies que l'impiété a faites à la France, d'y ranimer la foi, et de resserrer plus que jamais les nœuds sacrés qui unissent tous les catholiques français avec l'Église de Rome. C'est pour rendre cette mission du
% Source : PDF original, vue 084, côté droit, page manuscrite 169.
père commun des fidèles plus efficace, que notre bon maître en excita dans tous les cœurs à son approche une ardeur toute nouvelle. Remercions-le principalement d'avoir inspiré ces sentiments à notre Congrégation.

\section*{Annonce de l'arrivée du pape}
% Mention marginale de la page manuscrite 169.
À la nouvelle de la prochaine arrivée du souverain pontife dans nos murs, la vieille cité des martyrs fut émue, non pas comme Hérode, de ce trouble que produit une mauvaise conscience ; mais de ce sentiment d'agitation, résultat inévitable d'une grande joie. Les vrais enfants de l'Église, les congréganistes surtout ne s'arrêtèrent guère aux motifs du voyage ; ce n'étoit pas là ce qui pouvoit les toucher ; mais comment n'eussent-ils pas été remplis de la plus douce allégresse en pensant que nous allions avoir le bonheur de posséder au milieu de nous le chef de l'Église, le représentant de J. C. sur la terre. Aussi se préparoient-ils à le recevoir comme un père également vénéré et chéri.

La Congrégation s'étoit occupée plusieurs jours avant l'arrivée de l'illustre voyageur, de chercher les moyens d'assurer à sa réception tout l'éclat qui pouvoit aider à la dilatation des sentiments renfermés dans nos cœurs. Elle eut bien désiré pouvoir prendre elle-même directement une part active à l'espèce de triomphe dont elle vouloit voir accompagner son entrée ; mais le mystère dont elle étoit enveloppée rendoit la chose impossible. Elle concevoit même peu d'espoir de pouvoir être présentée, ce qui dans tous les cas ne pouvoit avoir lieu que d'une manière tout à fait instantanée. Cependant elle chargea M. Mas son secrétaire de préparer un discours. Les congréganistes firent en même temps des démarches officieuses auprès de l'autorité ecclésiastique pour obtenir qu'on mît le plus grand appareil à l'entrée du pape ; mais ne pouvant
% Source : PDF original, vue 085, côté gauche, page manuscrite 170.
pas faire en cette qualité des démarches officielles, ils les firent comme membres des confréries du St Sacrement ; ces sociétés avoient une existence publique et avouée ; il y avoit des congréganistes dans presque tous les bureaux de ces confréries ; il leur fut donc facile de déterminer ces bureaux à solliciter des ordres de la part de l'archevêché. Jusqu'au dernier moment, toutes ces démarches réitérées demeurèrent à peu près sans résultat. Le cardinal étoit absent, et ne devoit arriver qu'avec le pape. Les vicaires généraux n'avoient rien décidé par eux-mêmes ; en sorte qu'on vit arriver le 10, 9bre jour où le pape devoit être à Lyon, sans qu'on pût savoir encore de quelle manière on le recevroit. Les bureaux des confréries étoient tous dans l'attente, la Congrégation plus encore qu'aucune d'eux. Sur le qui-vive, les excitant par sous-main.

\section*{Démarches pour la réception du pape}
% Mention marginale de la page manuscrite 170 ; placée avant le dernier préparatif.
Sur les huit heures du matin, le bureau de la confrérie du St Sacrement de St Pierre s'étoit assemblé pour savoir si des ordres étoient parvenus. Lorsqu'il apprit qu'il n'y en avoit point encore, il se détermina à envoyer à l'archevêché trois députés pour connoître enfin quelle part la confrérie pourroit prendre à la réception. Sur ces trois députés, deux étoient congréganistes ; ils arrivent à l'archevêché et sont reçus par M. Jauffret, premier vicaire général, dans son cabinet. « Monsieur, disent-ils, on annonce l'arrivée du pape pour ce soir, nous venons vous prier de nous dire si les confrères du St Sacrement ont quelque chose à faire pour sa réception. — Mais je ne vois pas trop ; il arrivera ce soir sur les quatre heures ; les autorités vont le recevoir aux portes de la ville ; il doit descendre de voiture à la porte de St Jean, nous serons là pour le recevoir. — Eh ! vous n'avez point donné d'ordre aux paroisses ? — J'ai fait
% Source : PDF original, vue 085, côté droit, page manuscrite 171.
prier MM. les curés de se rendre à la cathédrale, ils se réuniront au chapitre ; je pense que pour éviter la confusion, on tiendra la porte fermée jusqu'après l'arrivée du pape. — Pardon, Monsieur le vicaire général, nous vous prions d'excuser la liberté que nous prenons ; mais aurez-vous le courage de repousser ainsi la foule empressée des fidèles ; elle ne pourra pas venir entourer le chef de l'Église, le père commun de tous ? — Je crois qu'il séjournera demain, et dans ce cas, nous espérons qu'il dira la messe à St Jean, le peuple pourra y assister. — Oui, mais en attendant votre réception aura lieu en cachette. Une véritable réception à la glace. Nos confrères avoient espoir qu'ils pourroient prendre part à cette réception. — Cela eût été sans doute très bien ; je conviens que cela va être un peu froid ; mais que voulez-vous que j'y fasse ? Maintenant il est trop tard pour y remédier. — Pas trop, tant que vous le croyez, Monsieur. Voulez-vous nous laisser faire ? Nous vous répondons de tout. Nous ne vous demandons qu'une seule chose, c'est d'être autorisés à faire dire de votre part à MM. les curés de se rendre cet après-midi à St Jean processionnellement avec leur clergé et leurs confréries du St Sacrement. » Les députés n'eurent pas plutôt arraché un demi consentement qu'ils courent à St Pierre où on les attendoit avec impatience.

Des lapins porteurs d'ordres\footnote{(Note du transcripteur) La lecture « lapins », confirmée par Timothée Gros, reste incertaine en raison de la singularité de la formule ; elle est conservée sans substitution.}, expédiés au nom de l'archevêché, partent aussitôt au pas de course pour toutes les paroisses de la ville ; partout on se hâte de convoquer les confréries. Les congréganistes avertis individuellement se chargent d'exciter de leur mieux le zèle des confrères ; à deux heures après midi des processions sortent simultanément de toutes les églises paroissiales. Les cloches sont mises toutes à la fois en grande volée. De nombreux
% Source : PDF original, vue 086, côté gauche, page manuscrite 172.
confrères en habits noirs et le flambeau à la main, marchent à la suite de la bannière de leur paroisse. Le chant du Benedictus retentit de toutes parts ; les magasins se ferment comme dans un jour de fête. Le peuple abandonne joyeusement ses ateliers, et court se presser aux environs de la primatiale ; déjà les confrères y sont arrivés, ils forment des deux côtés de la grande nef, une triple haie, illuminée par l'éclat de leurs flambeaux ; un nombreux clergé est réuni dans le chœur. Cette espèce d'agitation que cause inévitablement un bonheur bien vivement attendu, remue fortement tous les cœurs ; enfin la grosse cloche annonce que le pape est entré à Lyon ; le clergé se dirige vers la porte de l'église pour le recevoir ; peu de moments après, il repasse entre nos rangs en précédant le St Père ; le vénérable pontife s'avance lentement entre ces lignes de confrères prosternés, et répand sur eux ses plus abondantes bénédictions. Après avoir adoré J. C. dans le sacrement de son amour pendant quelques instants, le saint père reçoit la bénédiction du St Sacrement qui est donnée par le cardinal Fesch. Il se dirige ensuite sur l'archevêché en repassant par la grande porte. À sa suite toute la masse des confrères s'ébranle. Elle vient l'entourer des mille feux de ses flambeaux, et le conduit ainsi jusqu'à la porte du palais à travers les flots du peuple qui se précipite à genoux sur son passage. L'enthousiasme est à son comble ; il se fait jour par des acclamations ; les cris de Vive le St Père, commencent à se faire entendre aux environs de l'archevêché ; bientôt ils retentissent dans tous les quartiers de la ville. Des illuminations nombreuses et spontanées témoignent d'autant mieux de la joie et du bon
% Source : PDF original, vue 086, côté droit, page manuscrite 173.
esprit de la cité, que les autorités ne les avoient provoquées, ni par leurs invitations, ni par leurs exemples ; car aucun lampion n'étoit venu dissiper l'obscurité habituelle des bâtiments publics. Les confrères du St Sacrement ne parviennent qu'avec grande peine à se rallier à leurs bannières respectives au milieu de cette multitude qui encombroit toutes les avenues de l'archevêché. Ils retournent à leurs paroisses en mêlant aux acclamations du peuple, le chant joyeux du Te Deum.

Le lendemain, Sa Sainteté offrit à St Jean le saint sacrifice. Une foule immense y assiste. Je n'ai pas besoin, je pense, d'ajouter que la Congrégation tout entière en faisoit partie. Notre société fut encore complètement représentée dans l'après-midi de la même journée sur la place de Bellecour, lorsque le pape voulut bien s'y montrer au peuple, et lui donner la grande bénédiction papale avec indulgence plénière, ainsi qu'il la donne à Rome, \textit{Urbi et orbi}\footnote{(Note du transcripteur) Littéralement : « À la ville et au monde ».}, aux grandes solennités. Les ruines, tristes souvenirs des fureurs et du délire de l'esprit révolutionnaire, subsistoient encore ; mais elles avoient disparu sous la foule épaisse de chrétiens dont elles étoient couvertes. Tout ce peuple prosterné paroissoit vivement apprécier son bonheur de pouvoir environner le chef de l'Église du témoignage de son amour et de sa vénération, et d'être en même temps comblé de ses bénédictions.

Dans la soirée quelques congréganistes parvinrent jusqu'auprès du St Père ; ils crurent avoir avec eux des chefs de la Congrégation ; mais par suite d'une méprise ils s'en trouvèrent séparés. L'un d'eux prit la parole,
% Source : PDF original, vue 087, côté gauche, page manuscrite 174.
et exprima comme il put, bien mieux encore avec son cœur qu'avec le langage timide qui sortoit de sa bouche, les sentiments que nous éprouvions tous. Le bureau de la Congrégation se présenta peu de temps après ; il avoit avec lui le secrétaire désigné pour remplir les fonctions d'orateur. Celui-ci avoit eu soin de se munir de son discours, mais malheureusement il étoit trop tard. Le St Père fatigué de la journée s'étoit retiré. Il fut tout à fait impossible d'être admis. Il fallut donc se résigner à ce désappointement et se borner à transmettre le discours à Sa Sainteté ; nous eussions été trop heureux ; Dieu nous réserva cette petite humiliation, et nous imposa ce sacrifice pour nous maintenir dans l'esprit d'abnégation et d'humilité.

Dès le lendemain matin, le pape fut obligé de continuer son voyage ; il eût bien désiré pouvoir s'arrêter plus longtemps mais Napoléon pressoit ; il falloit qu'en toute hâte il se rendît à Paris. La réception qu'on venoit de lui faire à Lyon produisit un grand effet. C'étoit la première ville française qui s'étoit trouvée sur son passage, elle avoit donné l'élan ; le triomphe commencé dans l'antique métropole des Gaules, continua et se succéda de ville en ville. Le feu de l'enthousiasme gagnoit de proche en proche. Les hommes les moins religieux ne pouvoient se défendre d'une vive émotion, en rendant hommage au premier des pasteurs de l'Église, établi par le Sauveur lui-même pour paître et les agneaux et les brebis ; malgré eux ils éprouvoient au fond du cœur, un secret désir de rentrer enfin dans ce bienheureux bercail. Cette première grâce étoit un des fruits
% Source : PDF original, vue 087, côté droit, page manuscrite 175.
de la présence du vicaire de J. C. au milieu de nous. Ce fut toujours de cette manière, accompagné de ces témoignages de la vénération des fidèles, que le saint père parvint jusqu'à Paris.

\section*{Le pape à Paris}
% Mention marginale de la page manuscrite 175.
Ce mouvement religieux se fit bien aussi sentir dans la capitale ; toutefois il faut en convenir, il y fut plus froid que partout ailleurs. Cette grande ville renferme sans doute des hommes éminemment vertueux, mais ils sont loin d'être en majorité. Nous avons eu le soin de faire connoître à la Congrégation de Paris, la manière dont le pape avoit été reçu à Lyon. M. de Noailles, dont le zèle saisissoit avec ardeur toutes les occasions de se rendre utile à la religion, auroit bien désiré voir le même élan se manifester à Paris. Il se fut estimé bien heureux s'il avoit pu inoculer les sentiments dont son cœur étoit pénétré, à cette multitude qui toute de feu pour les plaisirs et les spectacles, est de glace lorsqu'il ne s'agit plus que des choses de la foi. Le saint père étoit venu visiter l'église de St Sulpice ; M. de Noailles se trouvoit sur le parvis ; dans le transport qui l'agite, il emploie toute l'énergie de ses poumons à grossir son organe, déjà en lui-même fort remarquable ; espérant enfin obtenir de l'écho, il faisoit retentir les airs de ses cris réitérés de : Vive le saint père. Mais hélas ! il reste seul, toujours seul. Aucune voix amie ne vient s'unir et répondre à la sienne ; tout à coup un individu en lui frappant sur l'épaule lui adresse ces paroles : courage, mon ami, courage, vous gagnez bien votre argent. Ce pauvre M. de Noailles étoit pris tout simplement pour un de ces clameurs à gage qu'on rencontre dans les théâtres de Paris.

Pendant le séjour du pape à Paris, une négociation qui
% Source : PDF original, vue 088, côté gauche, page manuscrite 176.
intéressoit bien vivement tous les catholiques lyonnais et en particulier la Congrégation, se traitoit dans notre ville.

\section*{Acquisition de Fourvières}
% Mention marginale de la page manuscrite 176.
Il étoit question de rentrer en possession de notre antique et vénérable église de Fourvières. Cette église avoit été vendue pendant la Révolution ; elle étoit tombée au pouvoir d'une femme qui l'avoit fait servir à l'exploitation d'une honteuse spéculation, elle recevoit à son profit les honoraires de messes que de crédules paysans venoient offrir, et faisoit ensuite acquitter ces messes au prix qu'il lui plaisoit de fixer, par des prêtres à ses ordres, schismatiques avec les schismatiques eux-mêmes, car M. Primat avoit rougi de cet odieux tripotage et avoit lancé contre eux une sentence d'interdiction, dont du reste ils n'avoient tenu aucun compte. Le cardinal s'étoit empressé de mettre fin à ce scandale en obtenant du gouvernement la clôture de l'église. La propriétaire se vit aussi privée du bénéfice simoniaque qu'elle en retiroit. Dans son désappointement il ne lui restoit plus d'autre parti à prendre que de vendre cette propriété. Mais elle vouloit en tirer un parti avantageux. Plusieurs établissements religieux particuliers en eussent fait volontiers l'acquisition ; l'archevêque et la masse des fidèles désiroient par-dessus tout que ce sanctuaire si cher à tous les Lyonnais n'appartînt qu'au diocèse. On entra en négociations ; la présence du souverain pontife en France en fit hâter la conclusion, on fut assez heureux pour que tout fût terminé à temps et il fut possible de faire coïncider avec le retour du pape, l'ouverture de ce sanctuaire de Marie.

Nous attendions le moment où nous devions posséder de
% Source : PDF original, vue 088, côté droit, page manuscrite 177.
nouveau ce pontife chéri avec d'autant plus d'empressement, que nous savions qu'il devoit séjourner plusieurs jours dans notre ville. L'entrain que sa présence avoit excité à son premier passage avoit été bien grand, et cependant les sentiments que nous éprouvions en voyant approcher le moment de son retour étoient bien plus vifs encore. En paroissant au milieu de nous pour la première fois, le pape en sa qualité de chef de l'Église avoit été entouré de toute notre vénération ; mais depuis que nous l'avions connu, il étoit devenu pour nous, s'il est possible, quelque chose de plus ; c'étoit un père, un père chéri, un père que nous avions déjà vu, dont nous avions pu admirer la douceur, la bonté, toutes les vertus. Nous allions donc le revoir. Ah ! nous nous disposions à l'accueillir comme une vieille connoissance, comme un ami. Ce ne sera plus seulement du respect que nous éprouverons en le voyant, ce sera de l'affection, ce sera de l'amour filial ; ce seront tous ces tendres sentiments qui remuent fortement le cœur des enfants bien nés, lorsqu'ils ont le bonheur, après une absence, de se retrouver dans les bras d'un excellent père.

Le retour du pape étoit annoncé pour le mardi de Pâques. Il avoit été jugé impossible de songer à nous rendre à St Jean comme nous l'avions fait à son arrivée. Par respect à Napoléon qui étoit fâché de voir les populations témoigner plus d'empressement pour la réception du pape que pour celle de l'empereur, le cardinal Fesch en avoit fait la défense formelle. Malgré ce mécompte, la Congrégation ne se tint pas pour battue ; elle usa de son influence sur les confréries du St Sacrement, pour engager les confrères
% Source : PDF original, vue 089, côté gauche, page manuscrite 178.
à se porter en grand nombre à la rencontre du pape. Le rendez-vous fut donné à la montée de Balmont, à quelque distance de la place de la Pyramide. On devoit s'y rendre isolément, mais en tenue de cérémonie, c'est-à-dire, en habit noir. Cet appel fut entendu ; quelques heures avant celle où le saint père devoit arriver on pouvoit apercevoir sur les deux quais de la Saône, des jeunes gens, par bandes de deux, de trois ou de quatre, se dirigeant tous vers le faubourg de Vaise. Cinq cents confrères parmi lesquels on auroit pu compter la Congrégation tout entière, se trouvèrent au lieu du rendez-vous. Ils se rangèrent en bataille sur la route. Le même orateur qui devoit prononcer le discours la première fois, fut encore chargé de haranguer le saint père ; pour nous c'étoit en qualité de secrétaire de la Congrégation, mais vis-à-vis des confrères du St Sacrement il fut désigné comme président des zélateurs de la confrérie de St Pierre. Il se plaça en avant de la cohorte entre deux confrères ; on avoit pris position au milieu de la montée sur un replat qui donnoit plus de facilité pour faire arrêter la voiture. Nous étions placés entre les autorités de Vaise qui attendoient le saint père au haut de la montée, et le préfet et les autres autorités de Lyon, qui se disposoient à le complimenter à la place de la Pyramide. Après une demi-heure d'attente, nous apercevons la voiture du pape escortée par la gendarmerie. Nous la voyons s'arrêter pour recevoir le compliment du maire de Vaise ; nous prenons à la hâte nos dernières dispositions ; voilà la voiture de nouveau en mouvement. Elle approche ; tous les cœurs sont dans l'attente ; le saint père est au milieu de ses enfants. Nous faisons signe au
% Source : PDF original, vue 089, côté droit, page manuscrite 179.
postillon d'arrêter, il hésite. Déjà les plus ardents d'entre nous, sautoient à la bride des chevaux. Le pape mit fin à ce débat en baissant lui-même la glace de la voiture et donnant l'ordre d'arrêter. L'orateur se présente à la portière, et adresse au vénérable pontife un compliment très court mais rempli d'une énergie toute chrétienne. Le pape témoigne sa satisfaction, et donne sa main à baiser aux trois personnes qui sont près de la portière et répand sa bénédiction sur tous les confrères, prosternés sur la route. Nous supplions Sa Sainteté de nous permettre d'environner sa voiture et de l'accompagner ainsi jusqu'à l'église cathédrale. Il nous répond avec un sourire dans lequel se peint avec vivacité toute l'expression de la bonté de son âme : \textit{Mult agréable à nous}. À l'instant nous formons un bataillon carré ; la voiture est au centre, et le cortège se met en marche aux cris mille fois répétés de Vive le St Père. Nous assistons ainsi sans quitter nos rangs, à la place de la Pyramide aux froides harangues officielles des autorités dans lesquelles Napoléon occupoit plus de place que le pontife vénéré auquel elles étoient adressées. On parcourt la grande rue du faubourg de Vaise ; puis on entre dans la ville. La population encombre les quais, la voiture s'avance lentement entre une double haie de chrétiens prosternés sur lesquels le St Père verse à pleines mains d'abondantes bénédictions. Tout fut au mieux jusque un peu après le rocher de Pierre Scise. Nous étions trop heureux, aussi le bon Dieu nous ménageoit-il une petite contrariété. Nos magistrats supportoient impatiemment notre présence autour de la voiture du pape. Napoléon venoit de séjourner plusieurs jours à Lyon et la population étoit demeurée très calme. L'enthousiasme qu'excitoit l'arrivée du chef de l'Église, offroit un point de comparaison qui devoit tout
% Source : PDF original, vue 090, côté gauche, page manuscrite 180.
naturellement paroître choquant à des hommes habitués à courber la tête devant les maîtres du pouvoir. Ils imaginèrent donc un moyen pour se débarrasser de nous et ils y réussirent ; ils font prendre le grand trot à l'escorte. La voiture du St Père suit le mouvement, malgré l'illustre voyageur qui en paroissoit affecté nous voilà forcés de demeurer en arrière. En nous retirant nous nous disions les uns aux autres, comme les disciples d'Emmaüs : ne nous sembloit-il pas que nos cœurs étoient brûlants, lorsque nous étions auprès de lui ? Deux congréganistes, MM. Eugène Roth et Loras restèrent dans le cortège ; ils se cramponnent à la voiture derrière laquelle ils se trouvoient placés, suivirent au pas de course et entrèrent à St Jean à la suite du souverain pontife. Ces deux amis, dans cette circonstance, représentoient bien réellement la Congrégation et cela de toute manière ; car c'est ainsi que bien ardente, dans la bonne comme dans la mauvaise fortune elle est restée constamment et invariablement attachée au chef de l'Église, et qu'elle espère avec la grâce d'en haut qui ne lui manquera pas davantage à l'avenir, lui demeurer toujours fidèle.

\section*{Séjour du pape}
% Mention marginale de la page manuscrite 180 ; placée au début des journées à Lyon.
Pendant les neuf beaux jours que le pape séjourna à Lyon, ce fut pour sa population une fête perpétuelle ; je n'ai pas besoin de vous répéter que les congréganistes prenoient une large part à toutes ces réjouissances. Partout où paroissoit le St Père ils se trouvoient toujours en grand nombre.

Le lendemain de son arrivée, le pape dit la messe à St Jean et communia de sa main tous les fidèles qui se présentèrent. Il y en eut un si grand nombre que quoique dans la prévision de cette affluence on eût consacré d'avance beaucoup d'hosties, elles manquèrent et il fallut avoir recours à un prêtre qui dit la
% Source : PDF original, vue 090, côté droit, page manuscrite 181.
messe pendant la communion pour en consacrer encore.

Le jour suivant le saint père, placé sur un trône élevé au fond de la basilique de St Jean, dans le même lieu où deux de ses prédécesseurs s'étoient assis il y a plusieurs siècles, admit tous les fidèles à lui baiser les pieds. Le soir il donna de nouveau comme il l'avoit fait à son premier passage, la bénédiction papale sur la place Bellecour. Comme la première fois, une foule immense remplissoit la vaste étendue de la place.

Pour fournir au peuple en masse une occasion de plus de pouvoir contempler les traits révérés du chef de l'Église, plus encore que pour lui faire connoître les rivages si riants de la Saône, on engagea le St Père à faire une promenade par eau jusqu'à l'Île Barbe. Une barque richement pavoisée avoit été préparée pour le recevoir. Il seroit difficile de retracer la magnificence de cette fête. Ces collines si richement décorées par la nature, aux pieds desquelles la Saône semble se complaire à promener doucement son onde calme et paisible, paroissoient en ce moment former sous la voûte du ciel un temple immense, dont la vaste étendue est remplie par la nombreuse population de la cité et des environs. La rivière elle-même a disparu sous cette multitude de gondoles élégamment ornées qui accompagnent la barque qui porte le chef de l'Église. Les spectateurs font retentir tous les échos d'alentour des acclamations réitérées de Vive le saint père. Mais dès que la nouvelle barque de Pierre, que leurs yeux suivent dans tous ses mouvements avec une pieuse avidité, paroit s'approcher du lieu où ils se trouvent placés, ils se jettent à genoux pour participer aux bénédictions que le saint père verse avec abondance sur son passage ; une seule pensée celle de rendre hommage
% Source : PDF original, vue 091, côté gauche, page manuscrite 182.
à J. C. fondateur et conservateur de la sainte Église, dans la personne de son vicaire, anime en ce moment tous les cœurs. Au retour de cette promenade, au moment où la barque touche au port, un arc-en-ciel tout rayonnant de lumière se forme dans les cieux, ce symbole de paix est l'arc de triomphe que le Seigneur a réservé au St Pontife représentant sur la terre le prince de la paix, et qui maintenant marchant à la suite et sur les traces de son divin maître, peut nous répéter après lui : apprenez de moi que je suis doux et humble de cœur.

\section*{Ouverture de Fourvières}
% Mention marginale de la page manuscrite 182.
Le 19 avril étoit le jour impatiemment attendu où l'église de Fourvières devoit enfin nous être rendue ; une solennité pareille appartenoit de droit à la Congrégation ; il étoit ici question du culte de Marie, notre mère ; c'étoit tout à fait pour les congréganistes une fête de famille ; aussi non seulement s'y rendirent-ils en grand nombre, mais ils en eurent en quelque sorte tous les honneurs de la journée. Sous le titre honorable de chevaliers de la Ste Vierge qui leur fut donné pour la circonstance, ils furent appelés pour la plupart à veiller à l'ordre de la cérémonie. Ils étoient tous en grande tenue, c'est-à-dire, en habit noir et gants blancs ; de plus pour marque de distinction, ils portoient tous au bras gauche, un ruban blanc passé au-dessus du coude. C'étoit la livrée de Marie.

Le souverain pontife entra dans l'église accompagné par une foule immense de peuple. Non seulement l'église, mais les cours et les chemins adjacents en étoient remplis. Le St Père offrit le saint sacrifice en présence de cette multitude, à l'autel qu'il rendoit au culte de Marie. Il recommanda sans doute à la mère de Dieu, cette bonne ville de Lyon qui depuis qu'elle a
% Source : PDF original, vue 091, côté droit, page manuscrite 183.
eu le bonheur de recevoir la foi, lui a été constamment dévouée. Après la messe du pape M. Courbon en dit une d'actions de grâces, pendant laquelle on chanta les litanies de la Ste Vierge. Le St Père entendit cette messe à genoux sur un prie-Dieu au milieu du sanctuaire. Lorsque la cérémonie fut terminée on fit entrer le pape dans une litière et MM. Caille, frères, et tous deux prêtres, se chargèrent en habit de chœur de ce précieux fardeau ; ils le transportèrent dans une petite maison qu'ils possédoient sur le plateau de Fourvières. Pour y aller on traversa les clos du voisinage, dont tous les propriétaires s'étoient empressés de tracer un chemin commode en ouvrant des brèches à leurs murs de clôture. Le peuple accompagnoit en foule et les chevaliers de la Ste Vierge formoient l'escorte. Dès que le pape se fut reposé un instant, on lui offrit son déjeuner habituel ; c'étoit une tasse de chocolat. On l'engagea ensuite à se rendre sur la terrasse pour y jouir d'un fort beau point de vue ; on découvre à la fois presque toute l'étendue de la ville, les gracieux détours du cours de la Saône, l'imposante majesté de celui du Rhône, les plaines immenses du Dauphiné, et pour couronner le tableau, les cimes altières des Alpes toujours couvertes de neiges.

Le moment est arrivé où après avoir replacé notre antique cité sous le patronage de la mère de Dieu, le successeur de St Pierre, du haut de cette montagne sanctifiée par le sang de nos premiers martyrs, va donner à la Rome de la France une bénédiction solennelle. Le saint père est placé au centre de la terrasse. Le clergé l'environne, tout autour les chevaliers de la Ste Vierge (c'est-à-dire la Congrégation) forment un demi-cercle ; derrière eux se presse une foule pieuse. On agite une bannière au-dessus
% Source : PDF original, vue 092, côté gauche, page manuscrite 184.
de la tête du vénérable pontife. Aussitôt les détonations de l'artillerie, se mêlent au son de toutes les cloches de la ville. Le vicaire de J. C. élève pieusement ses yeux et les mains vers le ciel, il étend ensuite sur la cité, ses mains remplies des bénédictions du Seigneur. Partout sur les quais, sur les ponts, sur les places publiques, et jusque sur les toits des maisons, on aperçoit le peuple prosterné. Chacun éprouve au fond du cœur un sentiment intime qui lui donne la conviction que la bénédiction que le saint père vient de donner à la terre est déjà confirmée dans le ciel.

\section*{Départ du pape}
% Mention marginale de la page manuscrite 184.
Le lendemain, fut un jour de deuil pour la cité ; nous éprouvions pour le père commun des fidèles, une tendresse toute filiale. Nous nous étions accoutumés à le voir au milieu de nous, et il fallut nous en séparer… Lorsqu'on se reporte sur les événements passés surtout s'il est question de scènes de sentiment, on est souvent tenté d'accuser l'historien d'exagération ; et cependant je puis bien vous l'assurer ; il n'y a point ici de pieuse fiction, mais seulement l'exacte vérité ; cette séparation fut pour toutes les âmes fidèles, l'occasion d'un sacrifice bien pénible. Au moment du départ du pape, un peuple nombreux environnoit l'archevêché ; on se précipitoit à genoux sur son passage, pour recevoir encore une fois sa bénédiction ; on crioit bien encore Vive le St Père mais d'une voix étouffée, tous les cœurs étoient serrés, tous les yeux se remplissoient de larmes.

Pontife vénérable et chéri, il n'est pas à douter qu'au milieu de toutes les amertumes dont votre cœur paternel a été abreuvé, le souvenir des hommages que vous avez reçus dans notre ville, de cette réception que vous avez bien voulu vous-même qualifier d'inexprimable, a apporté quelques consolations aux
% Source : PDF original, vue 092, côté droit, page manuscrite 185.
douleurs de ce cœur déchiré. Nous ne nous rappelons pas sans attendrissement que jusqu'à votre mort les Lyonnais ont toujours été vos enfants chéris, que ceux d'entre eux qui ont eu le bonheur de visiter la capitale du monde chrétien ont reçu de votre sainteté l'accueil le plus paternel, qu'elle a daigné plusieurs fois louer la Congrégation dans la personne de ses membres, et enfin qu'elle a ouvert pour elle les trésors de l'Église pour l'enrichir d'étonnantes indulgences. Maintenant très St Père, vous êtes en possession de ce beau ciel où nous avons la ferme confiance que vos vertus vous ont placé. Continuez, nous vous en supplions, à être notre protecteur et notre père, priez pour la cité de Pothin et des Irénée, afin que toujours unie à la chaire de St Pierre, elle conserve la foi catholique jusqu'à la fin des siècles ; priez, nous vous en conjurons plus vivement encore, priez pour cette Congrégation qui nous est si chère. Souvenez-vous combien elle étoit heureuse lorsqu'elle a pu contribuer à faire honorer la religion de J. C. par les hommages qu'elle s'efforçoit de rendre et de faire rendre au chef de l'Église ; obtenez-lui donc la grâce de persévérer dans sa vocation et de conserver à jamais, cet esprit de zèle et de charité qui doit être le sien ; que ses membres soient toujours entièrement dévoués à la cause de la foi ; qu'ils procurent autant qu'ils pourront dans ce monde la gloire de notre Dieu, afin qu'un jour dans le ciel ils soient jugés dignes de vous environner encore comme ils étoient si heureux de vous environner sur cette terre, et qu'unis avec vous et avec tous les saints, ils soient admis à chanter pendant l'éternité le cantique d'actions de grâces, au pied du trône de l'Agneau par lequel nous avons été sauvés.

Peu de jours après le départ du souverain pontife une relation de
% Source : PDF original, vue 093, côté gauche, page manuscrite 186, avant le chapitre septième.
son séjour, rédigée par un congréganiste fut lue en assemblée générale. Elle étoit destinée à être envoyée aux Congrégations associées, mais la Congrégation pensa que la publication de cette relation pourroit être utile ; elle ordonna donc qu'elle seroit imprimée. On y joignit les discours prononcés dans la circonstance, et notamment les deux discours adressés au pape au nom de la Congrégation, sous la simple indication de discours des jeunes gens de Lyon. Le tout forma une petite brochure qu'on intitula Lettre d'un jeune Lyonnais à un de ses amis, sur le passage du pape à Lyon. Cette brochure fut vendue au profit de la section des aumônes.
